\documentclass[conference]{IEEEtran}
\IEEEoverridecommandlockouts
% The preceding line is only needed to identify funding in the first footnote. If that is unneeded, please comment it out.
\usepackage{cite}
\usepackage{amsmath,amssymb,amsfonts}
\usepackage{algorithmic}
\usepackage{graphicx}
\usepackage{textcomp}
\usepackage{xcolor}
\def\BibTeX{{\rm B\kern-.05em{\sc i\kern-.025em b}\kern-.08em
    T\kern-.1667em\lower.7ex\hbox{E}\kern-.125emX}}
\begin{document}

\title{
\thanks{Identify applicable funding agency here. If none, delete this.}
}

\author{\IEEEauthorblockN{1\textsuperscript{st} Given Name Surname}
\IEEEauthorblockA{\textit{dept. name of organization (of Aff.)} \\
\textit{name of organization (of Aff.)}\\
City, Country \\
email address or ORCID}
\and
\IEEEauthorblockN{2\textsuperscript{nd} Given Name Surname}
\IEEEauthorblockA{\textit{dept. name of organization (of Aff.)} \\
\textit{name of organization (of Aff.)}\\
City, Country \\
email address or ORCID}}

\maketitle

\begin{abstract}
This survey provides a comprehensive overview of the current state of research on Embodied AI with Large Language Models (LLMs), highlighting the key themes, findings, and future directions in this rapidly evolving field. The integration of LLMs with Embodied AI has significant implications for decision-making, navigation, exploration, and question-answering tasks, and recent studies have demonstrated the potential of LLMs to improve performance in these areas. The survey examines the background and foundations of LLMs, their applications in Embodied AI, technical developments in architectures and training methods, and challenges and limitations associated with hallucinations, bias, and knowledge boundaries.

The analysis also explores emerging trends in multimodal LLMs, egocentric perception and understanding, human-LLM collaboration, and open challenges and future directions in LLM research. A critical examination of existing benchmarks and evaluation metrics for LLMs in Embodied AI reveals the need for comprehensive and diverse assessments. The survey highlights the role of specialized hardware, efficient training methods, and real-world deployments in advancing the field, with a focus on applications such as navigation, object manipulation, and environmental perception.

This survey contributes to the understanding of what makes Embodied LLMs effective, providing insights into their strengths, limitations, and potential applications. The scope of this research encompasses various domains, including natural language processing, robotics, and human-computer interaction. By synthesizing recent findings and identifying future directions, this survey aims to inform and guide researchers, practitioners, and policymakers in the development and deployment of Embodied AI systems with LLMs. Ultimately, this research has significant implications for the advancement of Embodied AI and its potential to transform various aspects of human life and society.
\end{abstract}

\begin{IEEEkeywords}
large language models, multimodal learning, natural language processing, machine learning, artificial intelligence
\end{IEEEkeywords}


\section{Introduction to Large Language Models}
Introduction to Large Language Models

The advent of large language models (LLMs) has revolutionized the field of natural language processing (NLP), enabling significant advancements in language understanding, generation, and decision-making tasks \cite{b1}. A comprehensive overview of LLMs is provided in \cite{b2}, which discusses their strengths, limitations, applications, and research directions. This survey highlights the potential of LLMs as powerful tools for various NLP tasks, while also emphasizing the need for critical evaluation and understanding of their capabilities and challenges.

The definition and scope of LLMs are critically examined in \cite{b3}, which provides a nuanced understanding of what constitutes an LLM. In contrast, other papers may have different understandings of LLMs, underscoring the importance of clarifying their definition and boundaries \cite{b4}. The evaluation metrics and methodologies used to assess LLMs also vary, with some papers proposing novel frameworks, such as the Embodied Agent Interface \cite{b5}, while others rely on more traditional approaches.

LLMs have been successfully applied in various domains, including language translation, text summarization, and dialogue systems \cite{b1}. Their effectiveness as embedding models is demonstrated in \cite{b4}, which surveys prompt designs, data-centric tuning methods, and advanced techniques for handling longer texts and multilingual data. The practical applications of LLMs are further explored in \cite{b6}, which offers a guide for working with LLMs, covering model usage, data influence, and downstream tasks.

A key theme emerging from the analysis of these papers is the importance of considering embodied contexts in LLM research \cite{b5}. The Embodied Agent Interface paper proposes a generalized interface for evaluating LLMs in embodied decision-making tasks, highlighting the need for more comprehensive evaluation frameworks to assess their performance. This underscores the significance of understanding how LLMs interact with their environment and make decisions in real-world scenarios.

Technical details, such as LLM architectures, prompt design, and data-centric tuning methods, are crucial in determining the performance of these models \cite{b4}. The choice of architecture, including GPT, BERT, and Mistral, can significantly impact the outcomes of LLM-based applications. Furthermore, the limitations and challenges associated with LLMs, such as hallucination errors and affordance errors, must be carefully evaluated to ensure their reliable deployment in real-world settings \cite{b5}.

In conclusion, this introduction to large language models highlights the current state of research in this field, including key findings, common themes, contrasting viewpoints, technical details, and limitations. By analyzing these developments, we can gain a deeper understanding of the potential and challenges of LLMs, informing future research directions and applications. The connections between LLMs and embodied decision-making tasks are particularly noteworthy, underscoring the need for more comprehensive evaluation frameworks and nuanced understandings of their capabilities and limitations \cite{b5}. As LLM research continues to evolve, it is essential to address these challenges and develop more effective methods for harnessing their potential in various domains.

\section{Background and Foundations of LLMs}
The development of Large Language Models (LLMs) has been a pivotal aspect of natural language processing, with significant implications for various applications. To understand what makes embodied LLMs effective, it is essential to delve into their background and foundations. This analysis draws upon five key papers that collectively provide an overview of the evaluation methods, semantic grounding, key claims, limitations, and career interests associated with LLMs.

The evaluation of LLMs is a multifaceted challenge, as highlighted by \cite{aSurveyOfUsefulLLMEvaluation}, which proposes a two-stage framework consisting of "core ability" and "agent" stages to assess their capabilities comprehensively. This approach underscores the need for refined methods to evaluate not just the linguistic prowess of LLMs but also their potential in real-world, embodied applications. The concept of semantic grounding is equally crucial, with \cite{understandingAISemanticGroundingInLargeLanguageModels} demonstrating that LLMs exhibit evidence of semantic grounding across functional, social, and causal dimensions. This finding suggests that LLMs are not merely stochastic parrots but possess a form of understanding, albeit one that requires further elucidation.

Defining LLMs and critically examining common claims about their properties is also vital, as emphasized by \cite{positionKeyClaimsInLLMResearchHaveALongTailOfFootnotes}. This critical examination process helps in identifying the strengths and limitations of LLMs, which are further elaborated upon by \cite{aPrimerOnLargeLanguageModelsAndTheirLimitations}. The latter highlights the applications and potential research directions for LLMs, underscoring their versatility while also cautioning against overlooking their limitations. Interestingly, \cite{theCareerInterestsOfLargeLanguageModels} reveals that LLMs tend to exhibit career interest inclinations towards social and artistic domains, which may not align with their demonstrated competences. This mismatch highlights the complexity of understanding LLMs' potential applications and their underlying motivations or biases.

A common theme across these papers is the need for more sophisticated evaluation methods and a deeper understanding of semantic grounding in LLMs. The recurring emphasis on evaluating LLMs beyond mere linguistic performance, as seen in \cite{aSurveyOfUsefulLLMEvaluation} and \cite{positionKeyClaimsInLLMResearchHaveALongTailOfFootnotes}, points to a broader recognition within the field of the importance of contextualizing LLM capabilities. Moreover, the discussion around semantic grounding, such as in \cite{understandingAISemanticGroundingInLargeLanguageModels}, suggests that there is a nuanced and multifaceted nature to how LLMs understand and interact with language.

The technical methodologies employed in these studies also offer insights into the complexities of analyzing LLMs. For instance, the application of core assumptions from theories of meaning in philosophy of mind and language to LLMs, as done by \cite{understandingAISemanticGroundingInLargeLanguageModels}, demonstrates an interdisciplinary approach to understanding these models. Similarly, the use of a general linear mixed model to analyze career interests, as seen in \cite{theCareerInterestsOfLargeLanguageModels}, showcases the diversity of analytical tools being applied to study LLMs.

Despite these advancements, challenges and limitations persist. Evaluating LLMs remains a daunting task due to their complexity and the multifaceted nature of their capabilities \cite{aSurveyOfUsefulLLMEvaluation}. Furthermore, fully comprehending semantic grounding and understanding in LLMs is an ongoing research endeavor \cite{understandingAISemanticGroundingInLargeLanguageModels}. The process of defining LLMs and critically examining claims about their properties is also fraught with challenges, requiring meticulous analysis and a critical perspective \cite{positionKeyClaimsInLLMResearchHaveALongTailOfFootnotes}.

In conclusion, the background and foundations of LLMs are characterized by a rich interplay of evaluation methodologies, semantic grounding, critical examinations of key claims, and explorations of their limitations and career interests. As research in this area continues to evolve, it is crucial to address the existing challenges and limitations, particularly in how we evaluate and understand the capabilities of LLMs. By doing so, we can harness the potential of embodied LLMs more effectively, paving the way for innovative applications that leverage their unique strengths while mitigating their weaknesses. Ultimately, a deeper understanding of what makes embodied LLMs work will depend on continued interdisciplinary research and a commitment to refining our methodologies and theories regarding these complex models.

\section{Applications of Large Language Models in Embodied AI}
The integration of Large Language Models (LLMs) with Embodied AI has emerged as a significant area of research, with far-reaching implications for decision-making, navigation, exploration, and question-answering tasks \cite{b1}. A key contribution in this domain is the proposed Embodied Agent Interface, which provides a generalized framework for evaluating LLMs in embodied decision-making tasks, offering valuable insights into their strengths and weaknesses \cite{b1}. This interface has been instrumental in highlighting the potential of LLMs to augment embodied intelligence systems with sophisticated environmental perception and decision-making support. Furthermore, the LLM-Brain framework demonstrates the potential of using LLMs as a robotic brain to unify egocentric memory and control, enabling effective exploration and question-answering tasks \cite{b2}. This approach underscores the importance of multimodal learning in embodied AI, where LLMs can be leveraged to capture complex spatial relationships and reason about the environment.

The use of multimodal LLMs has been a common theme in recent research, with applications in developing Generalist Embodied Agents (GEAs) that can generalize across diverse domains, including Embodied AI, Games, UI Control, and Planning \cite{b4}. This development is significant, as it enables the creation of agents that can adapt to new environments and tasks without requiring extensive retraining. Moreover, studies have shown that LLMs can capture implicit human intuitions about spatial building blocks of language, even without embodied experiences, as demonstrated in the reproduction of psycholinguistic experiments \cite{b3}. This finding has important implications for our understanding of the capabilities and limitations of LLMs in embodied contexts.

A survey of advances in embodied navigation using LLMs highlights the potential of these models to augment embodied intelligence systems with sophisticated environmental perception and decision-making support \cite{b5}. This research underscores the importance of evaluating LLMs in embodied contexts to understand their strengths and weaknesses, as well as the need for developing more generalizable frameworks for embodied AI. While some papers focus on the development of specialized frameworks for embodied AI, others propose more generalizable approaches, such as the LLM-Brain framework \cite{b2}. The use of supervised learning and online reinforcement learning (RL) to train GEAs is also contrasted with other approaches that rely on pre-trained models or benchmark-specific training \cite{b4}.

The technical details of these approaches vary, but most involve a combination of natural language processing, vision-language models, and reinforcement learning. For instance, the Embodied Agent Interface uses a combination of goal interpretation, subgoal decomposition, action sequencing, and transition modeling to evaluate LLMs in embodied decision-making tasks \cite{b1}. In contrast, the LLM-Brain framework employs a zero-shot learning approach, utilizing multimodal language models for robotic tasks and closed-loop multi-round dialogues for perception, planning, control, and memory \cite{b2}.

Despite these advances, there are several limitations and challenges associated with the integration of LLMs with Embodied AI. The evaluation of LLMs in embodied contexts is often limited by the availability of suitable datasets and benchmarks \cite{b1}. Furthermore, the development of generalizable frameworks for embodied AI requires significant advances in areas like multimodal learning and reinforcement learning \cite{b2}. The capture of human-like intuitions and spatial reasoning without explicit embodiment is also a challenging task, requiring further research into the capabilities and limitations of LLMs \cite{b3}. Finally, the development of GEAs faces challenges related to data quality, algorithmic choices, and the need for large-scale interactive simulators \cite{b4}.

In conclusion, the integration of LLMs with Embodied AI is a rapidly evolving field, with significant potential for advancing decision-making, navigation, exploration, and question-answering tasks. While there are several challenges associated with this research, the development of multimodal LLMs, GEAs, and more generalizable frameworks for embodied AI holds great promise for creating more sophisticated and adaptive embodied agents. Future research should focus on addressing the limitations and challenges associated with this field, including the need for more comprehensive datasets, benchmarks, and evaluation metrics. Ultimately, the successful integration of LLMs with Embodied AI has the potential to revolutionize a wide range of applications, from robotics and autonomous systems to human-computer interaction and cognitive architectures \cite{b5}.

\section{Technical Developments: Architectures and Training Methods for LLMs}
The development of large language models (LLMs) has been a rapidly evolving field, with significant advancements in architectures and training methods. A recent study, "Comprehensive Reassessment of Large-Scale Evaluation Outcomes in LLMs" \cite{b1}, challenges prevailing findings on emergent abilities and training types/architectures in LLMs, providing new perspectives on their characteristics and developmental trajectories. This reassessment highlights the importance of evaluation and assessment in understanding and improving LLMs, a theme that is echoed in other studies such as "A Survey of Useful LLM Evaluation" \cite{b3}, which proposes a two-stage framework for evaluating LLMs, focusing on core ability and agent applications.

The need for more nuanced and comprehensive approaches to evaluating LLMs is emphasized in the literature, including the consideration of multiple factors and sources of feedback. For instance, "Demystifying Issues, Causes and Solutions in LLM Open-Source Projects" \cite{b2} identifies common issues faced by practitioners, such as model issues, configuration problems, and connection difficulties, and potential solutions, including optimizing models. This study uses an empirical approach, collecting and analyzing data from 15 open-source projects to identify challenges and solutions, whereas "Comprehensive Reassessment" employs statistical techniques, including ANOVA, Tukey HSD tests, GAMM, and clustering methods \cite{b1}. The diversity of approaches highlights the complexity of LLM development and the need for multifaceted evaluation strategies.

Embodied reasoning and planning are emerging as key areas of research in LLMs, with a focus on leveraging environment feedback and language models to improve performance. "Inner Monologue: Embodied Reasoning through Planning with Language Models" \cite{b4} demonstrates the effectiveness of using environment feedback to improve high-level instruction completion in embodied contexts, proposing a novel approach that combines language models with planning algorithms. This study underscores the potential of embodied LLMs to enhance task execution and decision-making capabilities, particularly in dynamic environments.

The technical details and methodologies employed in these studies are noteworthy, as they reflect the evolving nature of LLM research. The use of advanced statistical techniques, such as those described in "Comprehensive Reassessment" \cite{b1}, and empirical approaches, like the one used in "Demystifying Issues" \cite{b2}, highlights the importance of rigorous evaluation and analysis in LLM development. Furthermore, the integration of environment feedback and planning algorithms, as seen in "Inner Monologue" \cite{b4}, opens up new avenues for research in embodied LLMs.

Despite these advancements, challenges and limitations persist. The papers highlight the need for more comprehensive and nuanced evaluation methods for LLMs, as well as the difficulties of working with complex and dynamic systems. Additionally, issues such as data quality, model interpretability, and scalability remain significant concerns \cite{b2}. The "Inner Monologue" study acknowledges the complexity of embodied reasoning and planning, emphasizing the need for further research in this area to fully realize the potential of LLMs \cite{b4}.

In conclusion, recent developments in architectures and training methods for LLMs have been significant, with a growing emphasis on evaluation, embodied reasoning, and planning. The studies discussed here demonstrate the diversity of approaches and techniques being explored, from statistical analysis to empirical studies and novel algorithmic frameworks. As LLM research continues to evolve, it is essential to address the challenges and limitations identified in these studies, ultimately paving the way for more effective, efficient, and embodied language models that can be integrated into real-world applications \cite{b1}, \cite{b2}, \cite{b3}, \cite{b4}.

\section{Challenges and Limitations of Current LLMs: Hallucinations, Bias, and Knowledge Boundaries}
The development and integration of Large Language Models (LLMs) in various applications have been accompanied by significant challenges and limitations, notably hallucinations, bias, and knowledge boundaries \cite{b1}. Hallucinations, which refer to the generation of content that is not based on any actual input or data, pose a substantial problem for the reliability and trustworthiness of LLMs \cite{b2}. Recent studies have explored the potential of hallucinations as a source of creativity rather than merely a limitation, suggesting that understanding and harnessing these phenomena could contribute to LLM applications by fostering innovative outputs \cite{b3}. However, the predominant view in the literature emphasizes the negative impact of hallucinations on the accuracy and reliability of LLMs, prompting research into mitigation strategies.

One approach to reducing hallucinations involves the incorporation of knowledge graphs into LLM architectures, which can help fill knowledge gaps and enhance reasoning accuracy by providing a structured representation of knowledge \cite{b4}. This technical solution underscores the importance of external knowledge augmentation in improving LLM performance. Additionally, there is a growing recognition of the need to redefine "hallucination" through psychological taxonomies that account for cognitive biases and other human psychological phenomena \cite{b5}. This perspective suggests that insights from human cognition can inform targeted solutions for addressing hallucinations in LLMs, highlighting the value of interdisciplinary approaches.

Surveys and analyses of recent efforts in detecting, explaining, and mitigating hallucinations in LLMs have presented comprehensive taxonomies of hallucination phenomena and discussed future research directions \cite{b6}. These studies emphasize the complexity and multifaceted nature of hallucinations, which make them challenging to define, detect, and mitigate. Furthermore, investigations into the perception of LLMs' knowledge boundaries using semi-open-ended question answering have introduced novel methods for discovering ambiguous answers that may indicate responses beyond an LLM's knowledge boundary \cite{b7}. This line of research underscores the significance of understanding and delineating what LLMs know and do not know, which is crucial for their effective and safe deployment.

A common theme across these studies is the acknowledgment of hallucinations as a significant challenge for LLMs, whether viewed as a potential source of creativity or primarily as a limitation to be overcome \cite{b8}. The need for external knowledge augmentation, such as through knowledge graphs, and the importance of understanding human cognition in addressing these challenges are also recurring patterns in the literature \cite{b9}. However, contrasting viewpoints exist, with some research focusing on the technical aspects of mitigating hallucinations and others adopting a more psychologically informed approach to understand and address these phenomena \cite{b10}.

The technical details and methodologies employed in these studies vary, ranging from the use of knowledge graphs and taxonomies of hallucinations to novel question-answering methods for perceiving knowledge boundaries \cite{b11}. Despite these efforts, significant limitations and challenges remain, including the complexity of hallucinations and the lack of standardized methods for evaluating the effectiveness of different mitigation strategies \cite{b12}. The development of innovative and interdisciplinary solutions is thus imperative for advancing the field and enhancing the reliability and performance of LLMs.

In conclusion, the challenges and limitations posed by hallucinations, bias, and knowledge boundaries in LLMs are complex and multifaceted, requiring comprehensive and innovative approaches to address them effectively \cite{b13}. Future research should prioritize the development of standardized evaluation methods, the exploration of interdisciplinary solutions combining insights from computer science, psychology, and other relevant fields, and the continued investigation into the potential of hallucinations as both a challenge and an opportunity for creativity and innovation in LLM applications \cite{b14}. By acknowledging these challenges and pursuing rigorous and interdisciplinary research, we can work towards creating more reliable, trustworthy, and capable LLMs that leverage their full potential while minimizing their limitations.

\section{Multimodal Large Language Models: Integrating Vision and Text for Embodied Cognition}
Multimodal large language models have emerged as a vital component of embodied cognition, enabling robots and agents to interact more effectively with their environments by integrating vision and text \cite{b1}. This integration has been shown to significantly improve the performance of robots in embodied tasks, such as object manipulation and navigation, by combining natural language instructions with visual perceptions \cite{b2}. The development of generalist embodied agents (GEAs) capable of grounding themselves across diverse domains through multi-embodiment action tokenizers has also been proposed, demonstrating strong generalization to unseen tasks \cite{b3}.

A common theme across these studies is the integration of different modalities, including text, vision, and action, to enhance the capabilities of embodied agents \cite{b4}. This multimodal approach allows for more effective interaction with complex environments and has been facilitated by advances in model architectures, such as multimodal GPT models like GPT-4V, and embodied language models like PaLM-E \cite{b5}. The introduction of comprehensive benchmarks, such as EmbodiedBench, has also highlighted the challenges in low-level manipulation tasks and the need for standardized evaluation platforms to accurately assess the performance of vision-driven embodied agents \cite{b6}.

Despite the progress made, several challenges remain, including adapting models to the complexity and variability of real-world environments \cite{b7}. The development of effective multimodal LLMs also requires access to large, diverse datasets that cover a wide range of tasks and scenarios \cite{b8}. Furthermore, finding the right balance between developing models that are specialized enough to perform well in specific tasks but also generalize well across different environments and tasks is an ongoing challenge \cite{b9}. Researchers have proposed various approaches to address these challenges, including the use of online reinforcement learning (RL) in interactive simulators to enhance adaptability and generalization \cite{b10}.

The choice between supervised learning and RL is a point of variation among studies, with some emphasizing the benefits of online RL for adapting to new environments and tasks \cite{b11}. Additionally, there is a contrast between efforts to develop generalist models capable of performing well across a wide range of tasks and those focusing on specialist models optimized for specific embodied tasks \cite{b12}. Open platforms for embodied agents are also seen as crucial for collective progress, allowing for more extensive experimentation and innovation \cite{b13}.

The integration of vision and text in LLMs has significant implications for the development of embodied cognition, enabling robots and agents to interact more effectively with their environments. As the field continues to evolve, addressing the challenges of adapting to real-world complexity, the need for diverse training data, and balancing specialization with generalization will be crucial for advancing the capabilities of embodied agents and their applications in real-world scenarios \cite{b14}. Ultimately, the development of multimodal LLMs has the potential to revolutionize the field of embodied cognition, enabling more effective human-robot interaction and paving the way for more sophisticated and autonomous robots \cite{b15}.

\section{Egocentric Perception and Understanding in Embodied AI with LLMs}
Egocentric perception and understanding play a pivotal role in embodied artificial intelligence (AI), particularly when integrated with large language models (LLMs). Recent studies have underscored the significance of unifying egocentric memory and control, as proposed by frameworks such as LLM-Brain, which utilizes LLMs as a robotic brain to achieve more integrated approaches in embodied AI \cite{b1}. This integration is crucial for enhancing dynamic scene understanding, as demonstrated by Embodied VideoAgent, which leverages both egocentric video and embodied sensory inputs to significantly improve performance in complex environments \cite{b2}. The importance of benchmarking egocentric video understanding capabilities is also highlighted through initiatives like VidEgoThink, which introduces a comprehensive benchmark for evaluating the effectiveness of models in first-person scenarios \cite{b3}.

A common theme across these studies is the emphasis on integrating multimodal information, including vision, language, and sensory inputs, to achieve enhanced performance in embodied AI tasks \cite{b4}. LLMs are recognized as crucial components in embodied AI systems, offering capabilities for perception, planning, and control when appropriately integrated with other modalities \cite{b5}. Furthermore, the challenges and opportunities in embodied learning and interaction, such as instruction following, exploration, and question answering, are frequently discussed, underscoring the need for comprehensive evaluation frameworks that can assess the capabilities of embodied AI systems in complex, dynamic environments \cite{b6}.

The technical details and methodologies employed in these studies also warrant attention. The use of multi-modal large models (MLMs) and world models (WMs) is highlighted for their potential to enhance perception, interaction, and reasoning in embodied AI \cite{b7}. Additionally, the development of custom benchmarks, such as VidEgoThink, underscores the importance of tailored evaluation metrics and benchmarks for assessing performance in specific embodied AI tasks \cite{b3}. Hybrid approaches combining symbolic and connectionist AI are also suggested as beneficial for achieving more robust and flexible embodied AI systems \cite{b8}.

Despite these advancements, several limitations and challenges remain. The current limitations of multi-modal large language models (MLLMs) in first-person scenarios are recognized as a significant challenge, necessitating further research to enhance their effectiveness in dynamic environments \cite{b9}. Moreover, the need for better integration of sensory inputs with LLMs is emphasized to achieve more seamless and effective interaction in embodied AI tasks \cite{b10}. Ensuring that embodied AI systems can operate safely and effectively, with appropriate levels of autonomy and human oversight, poses another significant challenge \cite{b11}.

In conclusion, the analysis of papers related to egocentric perception and understanding in embodied AI with LLMs highlights key themes and developments in this field. The integration of multimodal information, the role of LLMs, and the need for comprehensive evaluation frameworks are among the prominent topics discussed. While challenges persist, the ongoing research efforts aimed at addressing these limitations will likely pave the way for significant advancements in embodied AI, enabling more effective and autonomous interaction in complex environments \cite{b12}. As the field continues to evolve, it is essential to consider the implications of these developments on the design and deployment of embodied AI systems, ensuring that they are both effective and safe for human interaction \cite{b13}.

\section{Human-LLM Collaboration: Replicating Crowdsourcing Pipelines and Enhancing Reasoning Abilities}
Human-LLM collaboration has emerged as a significant area of research, with a focus on replicating crowdsourcing pipelines and enhancing reasoning abilities. Studies have shown that large language models (LLMs) can replicate human-like behavior in crowdsourcing tasks, but their success is variable and influenced by factors such as requester understanding of LLM capabilities and task requirements \cite{b1}. Replicating crowdsourcing pipelines offers a valuable platform to investigate LLM strengths and weaknesses on different tasks and their potential in complex tasks \cite{b1}. Moreover, LLMs can improve the effectiveness of crowdsourcing workflows, but more research is needed to understand how to evaluate and design these workflows \cite{b2}.

Techniques from crowdsourcing workflows can be adapted to develop LLM chains, which enable complex tasks by decomposing work into a sequence of subtasks \cite{b3}. This approach has been shown to be effective in achieving more efficient and effective outcomes. Furthermore, embodied LLMs can reason over sources of feedback provided through natural language, allowing them to form an inner monologue that improves planning and decision-making in robotic control scenarios \cite{b4}. Organized team structures can be imposed on LLM agents to foster cooperation and reduce communication costs, with designated leadership having a positive impact on team efficiency \cite{b5}.

A common theme among these studies is the potential of LLMs to replicate human-like behavior in various tasks, including crowdsourcing and embodied reasoning. Understanding LLM capabilities and limitations is crucial in designing effective workflows and systems. The need for more research on evaluating and designing LLM-based systems, particularly in complex and dynamic environments, is also a recurring theme. Feedback and interaction play a significant role in improving LLM performance and decision-making, highlighting the importance of developing mechanisms for effective communication between humans and LLMs.

While there are similarities among the studies, contrasting viewpoints and approaches are also evident. For instance, \cite{b1} focuses on replicating crowdsourcing pipelines with LLMs, whereas \cite{b2} takes a more general approach to exploring the potential of LLMs in crowdsourcing workflows. \cite{b3} adapts techniques from crowdsourcing workflows to develop LLM chains, whereas \cite{b4} focuses on embodied reasoning and planning with LLMs. These differences in approach highlight the diversity of research in human-LLM collaboration and the need for further investigation.

The technical details and methodologies employed in these studies vary, including crowdsourcing platforms and workflows \cite{b1}, \cite{b2}, \cite{b3}, LLM chain development and evaluation \cite{b3}, embodied reasoning and planning with LLMs \cite{b4}, and organized team structures and leadership models for LLM agents \cite{b5}. Experimental evaluations and case studies are used to demonstrate the effectiveness of proposed approaches, providing valuable insights into the potential applications and limitations of human-LLM collaboration.

Despite the progress made in this area, several limitations and challenges remain. The need for more research on evaluating and designing LLM-based systems is a significant challenge, particularly in complex and dynamic environments \cite{b2}. The potential risks and challenges of relying on LLMs in such environments must also be addressed \cite{b1}. Additionally, issues such as information redundancy and confusion in multi-agent cooperation with LLMs need to be resolved \cite{b5}. Developing more effective organizational structures and leadership models for LLM agents is also essential \cite{b5}. Ultimately, scaling up LLM-based systems to real-world applications and scenarios remains a significant challenge, requiring further research and development.

In conclusion, human-LLM collaboration has the potential to revolutionize various areas, including crowdsourcing, embodied reasoning, and multi-agent cooperation. The studies analyzed in this section highlight key themes and developments, including the importance of understanding LLM capabilities and limitations, the need for effective communication mechanisms, and the potential benefits of combining human and LLM strengths. As research in this area continues to evolve, it is essential to address the limitations and challenges identified, ultimately leading to more efficient and effective human-LLM collaboration systems.

\section{Open Challenges and Future Directions in LLM Research: Ethics, Explainability, and Generalizability}
The development of embodied systems integrated with Large Language Models (LLMs) has sparked significant interest in recent years, with potential applications spanning various domains, including human-computer interaction, robotics, and natural language processing \cite{b1}. However, as LLMs become increasingly sophisticated and pervasive, several open challenges and future directions emerge, particularly concerning ethics, explainability, and generalizability. A critical examination of the current state of LLM research reveals that ethical considerations are paramount, with a growing recognition of the need for responsible and ethically aligned language models \cite{b2}. The importance of ethics in LLM research is underscored by the potential risks and consequences of deploying these models without proper consideration of their ethical implications, including issues related to bias, fairness, transparency, and accountability \cite{b3}.

The analysis of relevant papers highlights the complexity of ethical concerns in LLM research, with longstanding issues such as copyright infringement, systematic bias, and data privacy being complemented by emerging problems like truthfulness and social norms \cite{b4}. In response to these challenges, researchers have proposed comprehensive guides to ethical research with LLMs, including best practices and toolkits to support ethical work, such as the LLM Ethics Whitepaper \cite{b5}. Furthermore, the introduction of corpora like EthiCon, which can be used to automate concern identification and provide insights into the ethical concerns of LLM researchers, represents a significant step towards addressing these challenges \cite{b6}. The development of frameworks for in-context ethical policies in LLMs, integrating moral dilemmas with moral principles, also offers a promising approach to infusing generic ethical reasoning capabilities into these models \cite{b7}.

Despite these advancements, the field of LLM research is not without its tensions and debates. While some researchers advocate for standardized ethical principles and guidelines, others propose more flexible and adaptive approaches to ethics, acknowledging the value pluralism and cultural diversity inherent in LLM development \cite{b8}. The use of natural language processing (NLP) techniques, such as those employed in the extraction and analysis of ethical concern statements, has been instrumental in shedding light on these complexities \cite{b9}. Moreover, the integration of machine learning models, like GPT-x, to test frameworks for in-context ethical policies underscores the interdisciplinary nature of LLM research \cite{b10}.

Looking ahead, it is clear that addressing ethical concerns in LLM research will require continued effort and collaboration from researchers, practitioners, and stakeholders. The limitations of current approaches, including the lack of standardized principles and guidelines, the need for more diverse and representative datasets, and the potential for bias in LLMs, must be acknowledged and addressed \cite{b11}. Ongoing monitoring and evaluation of LLMs will be essential to ensure that they are aligned with ethical standards and societal values, and that their development is guided by a commitment to responsible innovation \cite{b12}. Ultimately, the future of embodied systems integrated with LLMs depends on our ability to navigate these challenges and develop models that are not only highly performant but also ethically sound, transparent, and generalizable. By prioritizing ethics, explainability, and generalizability in LLM research, we can unlock the full potential of these technologies while minimizing their risks and ensuring that they serve the greater good \cite{b13}.

\section{Survey of Existing Benchmarks and Evaluation Metrics for LLMs in Embodied AI}
The survey of existing benchmarks and evaluation metrics for Large Language Models (LLMs) in embodied AI reveals significant contributions to the understanding and assessment of these models in decision-making tasks. A critical analysis of relevant papers, including \cite{b1}, \cite{b2}, and \cite{b3}, highlights several key findings that underscore the importance of comprehensive and diverse benchmarks. For instance, the proposal of a generalized interface for benchmarking LLMs in embodied decision-making tasks, as seen in \cite{b1}, emphasizes the need for fine-grained metrics beyond traditional success rates. This is further supported by \cite{b2}, which introduces an interactive evaluation benchmark featuring 328 distinct tasks across varied 3D scenes, thereby revealing significant shortcomings of state-of-the-art models compared to human performance on embodied tasks.

The introduction of comprehensive benchmarks such as EmbodiedBench \cite{b4}, with its 1,128 testing tasks, evaluates essential agent capabilities like commonsense reasoning and long-term planning. It is noteworthy that while multimodal LLMs (MLLMs) excel in high-level tasks, they struggle with low-level manipulation, indicating a gap in current model capabilities. The analysis of inadequacies in LLM benchmarks by \cite{b5} points out limitations such as biases, difficulties in measuring genuine reasoning, and the need for standardized methodologies and ethical guidelines. Furthermore, \cite{b6} provides an overview of advancements in embodied AI, emphasizing the role of multi-modal large models (MLMs) and world models (WMs), and discusses challenges and future directions, thus offering a broad perspective on the field.

A common theme across these papers is the recognition of the importance of comprehensive and diverse benchmarks for evaluating LLMs in embodied AI. There is a consensus on the need for benchmarks that go beyond simple success rates, incorporating fine-grained metrics to assess various aspects of performance \cite{b1}, \cite{b2}. The diversity and complexity of tasks are also highlighted as crucial, with an emphasis on testing LLMs with tasks that mimic real-world scenarios, including both high-level reasoning and low-level manipulation \cite{b3}, \cite{b4}. Multimodal interaction is another significant aspect, reflecting the need for models that can effectively perceive, interact, and reason across different modalities \cite{b5}. Ethical considerations and standardization are also pointed out as essential for ensuring fairness, transparency, and safety in LLM evaluation \cite{b6}.

While there is a shared goal of improving LLMs in embodied AI, contrasting viewpoints emerge regarding the approach. For example, the critique of current benchmarks suggests a shift from static to dynamic behavioral profiling \cite{b7}, reflecting a desire for more adaptive and responsive evaluation methodologies. Additionally, some papers focus primarily on technical aspects \cite{b1}, \cite{b2}, while others emphasize ethical considerations and the need for regulatory guidelines \cite{b5}, \cite{b6}. Technically, these contributions involve novel benchmarking frameworks such as EmbodiedEval and EmbodiedBench, which provide structured ways to evaluate LLMs across a wide range of tasks. There is also an emphasis on developing metrics that can accurately capture nuanced performance aspects, such as reasoning capabilities and adaptability \cite{b3}, \cite{b4}.

Despite these advancements, several limitations and challenges are identified, including the inadequacies of current benchmarks, the need for standardization, ethical and regulatory challenges, and technical hurdles \cite{b5}, \cite{b6}. The absence of standardized methodologies hinders comparison and progression in the field, while ensuring that LLMs are developed and used responsibly poses ethical and regulatory challenges. Improving performance, especially in low-level manipulation tasks, and developing models that can effectively interact in dynamic environments remain significant technical hurdles.

This analysis highlights the evolving landscape of LLM evaluation in embodied AI, emphasizing the need for comprehensive, diverse, and ethically sound benchmarks. It underscores the challenges ahead but also points to the potential for innovation and advancement through collaborative efforts and a commitment to addressing both technical and ethical considerations. As the field continues to evolve, it is crucial to prioritize the development of standardized, dynamic, and inclusive evaluation frameworks that can accurately assess the capabilities of LLMs in embodied AI and guide future research towards more effective, safe, and responsible solutions \cite{b1}, \cite{b6}.

\section{Emerging Trends: Specialized Hardware, Efficient Training Methods, and Real-World Deployments}
The field of embodied intelligence with Large Language Models (LLMs) is witnessing significant advancements, driven by emerging trends in specialized hardware, efficient training methods, and real-world deployments. A key finding in this area is the effectiveness of LLMs in navigation tasks that require a deep understanding of the environment and quick decision-making \cite{b2}. This has led to the development of novel architectures, such as Learnable Latent Codes as Bridges (LCB), which aim to overcome the limitations of using language as an interface layer between high-level task planners and low-level policies \cite{b3}. The creation of open platforms like LEGENT has also facilitated the integration of LLMs and Large Multimodal Models (LMMs) into embodied agents, promoting collective progress in this field \cite{b4}.

The importance of efficient training methods and specialized hardware to support the deployment of LLMs in real-world applications is a recurring theme across research papers \cite{b1, b5}. This emphasizes the need for innovative solutions that can improve performance and efficiency without altering core decoding mechanisms. In line with this, advances in LLM serving systems have focused on system-level enhancements, highlighting the significance of interdisciplinary approaches that combine hardware and software expertise \cite{b1}. Furthermore, the use of multimodal models and fusion of language and vision capabilities is a growing trend in embodied intelligence research, enabling more effective and flexible interaction with environments \cite{b2, b4}.

While some researchers focus on developing novel architectures to overcome limitations of language-based interfaces \cite{b3}, others propose using open platforms to facilitate the integration of LLMs and LMMs into embodied agents \cite{b4}. This diversity in approaches underscores the complexity and multifaceted nature of embodied intelligence, requiring a range of solutions tailored to different applications and contexts. Technical details, such as the use of learnable latent codes by LCB to enable flexible communication of goals without language limitations \cite{b3}, and LEGENT's dual approach featuring interactive 3D environments and sophisticated data generation pipelines \cite{b4}, demonstrate the sophistication and nuance of current research.

However, challenges persist, including domain shift and catastrophic forgetting when fine-tuning LLMs on embodied data \cite{b3}, as well as ethical considerations related to the deployment of LLMs in real-world scenarios \cite{b5}. Scalability and efficiency of LLM serving systems also remain critical for widespread adoption, necessitating further research and development \cite{b1}. Despite these challenges, the emerging trends in specialized hardware, efficient training methods, and real-world deployments signal a promising future for embodied intelligence with LLMs. As research continues to address the technical, ethical, and practical aspects of this field, we can expect significant advancements in the capability and applicability of embodied agents, leading to innovative solutions across various domains.

\section{Conclusion and Outlook: The Role of Large Language Models in Shaping the Future of Embodied AI}
The integration of Large Language Models (LLMs) with Embodied AI has emerged as a significant area of research, with potential applications in various domains. Recent papers have highlighted the role of LLMs in shaping the future of Embodied AI, emphasizing their potential to improve decision-making, environmental perception, and navigation \cite{b1}, \cite{b3}. A proposed embodied agent interface provides a generalized framework for evaluating LLMs in embodied decision-making tasks, allowing for a comprehensive assessment of their performance and identification of strengths and weaknesses \cite{b1}. Moreover, studies have shown that LLMs can capture implicit human intuitions about spatial building blocks of language, despite their non-embodied nature, revealing adaptability without a tangible connection to embodied experiences \cite{b2}.

The development of generalist embodied agents is also a key area of research, with multimodal large language models (MLLMs) being adapted to enable grounding across diverse domains through a multi-embodiment action tokenizer \cite{b4}. This highlights the importance of multimodal learning, allowing agents to process and integrate information from various sources, such as language, vision, and sensorimotor experiences. Furthermore, embodied AI is proposed as the next fundamental step in the pursuit of Artificial General Intelligence, emphasizing the importance of creating agents capable of seamless communication, collaboration, and coexistence with humans and other intelligent entities \cite{b5}.

The integration of LLMs with Embodied AI has been explored through various approaches, including task-specific models \cite{b1} and generalist agents \cite{b4}. While non-embodied LLMs can still capture spatial schema intuitions \cite{b2}, the importance of embodied experiences for Artificial General Intelligence cannot be overstated \cite{b5}. Technical details, such as the use of supervised learning and reinforcement learning to evaluate LLMs in embodied decision-making tasks \cite{b1}, and the employment of multimodal learning techniques, including online RL and cross-domain data training \cite{b4}, have been crucial in advancing the field.

However, several limitations and challenges need to be addressed, including the need for advanced hardware to enable agents to interact with their environment and process complex sensory information. Additionally, the formulation of a novel AI learning theory, capable of integrating multimodal information and supporting the development of generalist agents, is essential \cite{b5}. Scalability and generalization also pose significant challenges, requiring further research to ensure the applicability of proposed models and approaches in real-world scenarios.

In conclusion, the integration of LLMs with Embodied AI has the potential to revolutionize various domains, from robotics to healthcare. As researchers continue to explore this intersection, it is essential to address the limitations and challenges outlined above, driving innovation in hardware, software, and theoretical frameworks. By doing so, we can unlock the full potential of Embodied AI, enabling the development of generalist agents capable of adapting to diverse domains and interacting seamlessly with humans and other intelligent entities \cite{b4}, \cite{b5}. Ultimately, this will pave the way for significant advancements in Artificial General Intelligence, transforming the way we interact with and understand the world around us.

\section{Conclusion}
This survey has comprehensively analyzed recent papers related to Embodied AI with Large Language Models (LLMs), providing insights into the current state of research in this field. The key findings suggest that LLMs can significantly improve their performance as a world model by extracting necessary information from visual data \cite{b1}. Furthermore, the concept of an inner monologue has been proposed, enabling LLMs to process and plan in robotic control scenarios by leveraging environment feedback \cite{b2}. The development of a generalized interface, such as the Embodied Agent Interface, has also been highlighted as a crucial factor in evaluating LLMs for embodied decision making, offering a comprehensive assessment of their performance \cite{b3}. Additionally, learnable latent codes have been identified as a potential bridge between LLMs and low-level policies, facilitating flexible communication of goals in task plans without language limitations \cite{b4}.

A common theme throughout the surveyed papers is the integration of LLMs with embodied intelligence, underscoring the potential benefits of combining these two areas. The utilization of visual information is also a prevalent trend, with several papers emphasizing its importance in improving the performance of LLMs in embodied AI applications. Moreover, the need for effective interfaces between LLMs and low-level policies or embodiment has been consistently highlighted, with various approaches proposed to address this challenge. While some papers advocate for language-based interfaces \cite{b2}, others suggest using latent code-based interfaces as a more effective alternative \cite{b4}. The evaluation methodologies employed in the surveyed papers also vary, ranging from benchmarking \cite{b3} to comparative performance analysis on specific tasks \cite{b4}.

The technical details and methodologies employed in the surveyed papers are diverse, with various LLM architectures, such as GPT-4V \cite{b4}, and embodiment platforms, including Minecraft \cite{b1} and robotic platforms \cite{b2}, being utilized. Latent code learning has also been employed to bridge the gap between LLMs and low-level policies \cite{b4}. However, despite these advancements, several limitations and challenges have been identified, including language limitations \cite{b4}, domain shift, and catastrophic forgetting \cite{b4}, as well as the need for effective evaluation methodologies \cite{b3}.

The implications of these findings are significant, suggesting that the integration of LLMs with embodied intelligence has the potential to revolutionize various applications, from robotics to computer vision. The development of effective interfaces and evaluation methodologies will be crucial in unlocking this potential. Future research should focus on addressing the identified limitations and challenges, exploring new approaches to integrate LLMs with embodied intelligence, and developing more comprehensive evaluation frameworks. By doing so, the field of Embodied AI with LLMs can continue to advance, leading to innovative solutions and applications that can transform the way we interact with and understand the world around us. Ultimately, this survey provides a foundation for understanding the current state of research in Embodied AI with Large Language Models, highlighting key themes, developments, and implications that will shape the future of this exciting and rapidly evolving field.

\end{document}